\section{Introduction}\label{sec:intro}

Prescriptive process management uses data and models to recommend or take decisions
in running business processes, such as which resource handles which case, which
activity to perform next, or whether a case needs an extra check. Reinforcement
learning (RL) is a natural fit when such decisions interact over time and their
consequences are delayed, and a growing body of work trains RL policies in simulation
models of business processes~\cite{aepn2023,gympn2025}. \todo{add RL-for-BPM
references (resource allocation, next-best-action, prescriptive monitoring)}

Real organisations, however, do not run one process. An insurer handles claims,
underwriting and complaints; a hospital runs an emergency department, outpatient
clinics and a pharmacy; a shared-service centre runs payables, receivables and payroll.
These processes typically have dedicated teams and run side by side, while
performance is measured on organisation-level indicators such as total throughput or
profit. Treating such an organisation as one monolithic RL problem has three
consequences:
\begin{enumerate}
  \item \emph{Diluted credit.} Each decision is credited with the reward noise of
    every process, including processes it cannot influence. No critic can remove this
    noise: it is the other processes' future randomness.
  \item \emph{All-or-nothing waiting.} Good resource allocation sometimes requires
    \emph{not} assigning available work, for example keeping a slow resource idle
    rather than giving it a job a faster resource will soon handle~\cite{linkumar1984}.
    In a monolithic model the only way to wait is to freeze the whole organisation.
  \item \emph{Cost that grows with size.} Every decision reasons over the whole
    organisation, so the cost of a decision grows with the organisation even when
    only one process is involved.
\end{enumerate}

Our starting observation is that the process model already contains the information
needed to avoid all three. In an Action-Evolution Petri Net (\aepn)~\cite{aepn2023},
two parts of an organisation can influence each other only through shared places,
such as a queue or a resource pool. The connected components of the net therefore
\emph{certify} which processes are independent. We use this independence certificate
consistently at three points of the RL pipeline:
\begin{itemize}
  \item \textbf{Credit}: Net-Factored GAE (\nfgae) credits each decision only with the
    rewards of its own component, on that component's own decision clock.
  \item \textbf{Decision timing}: component-scoped postponement lets one process wait
    for a better resource while the others keep working.
  \item \textbf{Observation}: component turns let each decision observe and encode only
    its own component.
\end{itemize}
Each use is exact where the model certifies independence and reduces to standard
proximal policy optimisation (PPO)~\cite{ppo2017} with generalised advantage
estimation (GAE)~\cite{gae2016} when the net has a single component.

\paragraph{Contributions.}
\begin{enumerate}
  \item The independence certificate for \aepn{}s and its three uses, with soundness
    conditions, unbiasedness, variance dominance over PPO and exact equivalence to PPO
    on a single process (Sect.~\ref{sec:method}).
  \item Component-scoped postponement, which lets an agent learn policies that
    deliberately leave a resource idle, per process and without coupling the processes; we show it is learnable only
    together with factored credit.
  \item An organisation-level benchmark: an insurer with a claims process
    (next-activity selection), an underwriting process (resource assignment) and a
    complaints process (rework and quality gate), with regions, a shared-team coupling
    variant and a variant in which waiting pays (Sect.~\ref{sec:setup}).
  \item An empirical study showing that standard PPO degrades as concurrency grows
    while the independence-aware pipeline does not, that it learns to wait where
    waiting pays and nowhere else, and that its cost grows linearly with the number of
    processes (Sect.~\ref{sec:results}).
\end{enumerate}
In short: on an insurer with six concurrent processes, joint PPO reaches 0.75 of the
heuristic's return after 40 epochs and factored credit 1.07, ahead on all 20 seeds and
at the heuristic's level after a quarter of the epochs; where leaving an underwriter idle
pays, only factored credit with component-scoped postponement learns to do so, beating
the waiting heuristic on 18 of 20 seeds with one region and 12 of 20 with two, while PPO
with the same option never does with two regions.
